% BHU PYQ -- Manually transcribed & error-corrected
% Paper: GLB-501: Physical and Structural Geology
% Exam: B.Sc. (Hons.) Semester V Examination, 2013-14

\documentclass[11pt,a4paper]{article}
\usepackage[a4paper,top=2.5cm,bottom=2.5cm,left=2.8cm,right=2.8cm,headheight=14pt]{geometry}
\usepackage{amsmath,amssymb}
\usepackage[T1]{fontenc}
\usepackage[utf8]{inputenc}
\usepackage{lmodern,microtype,fancyhdr,enumitem,hyperref,lastpage,parskip}

\hypersetup{colorlinks=true,urlcolor=black,linkcolor=black,
  pdftitle={GLB-501 Physical and Structural Geology -- BHU B.Sc. Sem V 2013-14}}
\pagestyle{fancy}\fancyhf{}
\renewcommand{\headrulewidth}{0.4pt}\renewcommand{\footrulewidth}{0.4pt}
\fancyhead[L]{\small\textit{Banaras Hindu University}}
\fancyhead[R]{\small\textit{GLB-501: Physical and Structural Geology}}
\fancyfoot[C]{\small Page \thepage\ of \pageref{LastPage}}

\newlist{parts}{enumerate}{2}
\setlist[parts,1]{label=(\alph*),leftmargin=*,itemsep=5pt,topsep=3pt}
\newlist{romanparts}{enumerate}{2}
\setlist[romanparts,1]{label=(\roman*),leftmargin=*,itemsep=5pt,topsep=3pt}
\newcommand{\pts}[1]{\hfill\textit{[#1]}}

\begin{document}
\begin{center}
  {\large\bfseries BANARAS HINDU UNIVERSITY}\\[2pt]
  {\normalsize B.Sc. (Hons.) --- Semester V Examination, 2013-14}\\[6pt]
  \rule{0.8\linewidth}{0.5pt}\\[6pt]
  {\large\bfseries Paper No. GLB-501: Physical and Structural Geology}\\[4pt]
  {\normalsize Time: 3 Hours \hfill Maximum Marks: 70}\\[2pt]
  {\small Ref. No.: BS/Sem V/257}
\end{center}
\rule{\linewidth}{0.4pt}
\smallskip
\noindent\textit{Instructions: Attempt any \textbf{five} questions in all, at least \textbf{two} questions each from Section A and B. Section C is compulsory. All questions carry equal marks (14 marks each).}
\bigskip

\bigskip\noindent{\large\bfseries SECTION --- A}
\rule{\linewidth}{0.4pt}\smallskip

\noindent\textbf{Question 1.} What is a fold? Give the genetic classification of folds and discuss the mechanism of their development. \pts{14}

\medskip\noindent\textbf{Question 2.} Describe the erosional and depositional features formed by the wind. \pts{14}

\medskip\noindent\textbf{Question 3.} Discuss the following: \pts{$2 \times 7 = 14$}
\begin{parts}
  \item Continental drift theory of Taylor
  \item Effect of faulting on folded outcrop
\end{parts}

\bigskip\noindent{\large\bfseries SECTION --- B}
\rule{\linewidth}{0.4pt}\smallskip

\noindent\textbf{Question 4.} Write short notes on the following: \pts{$4 \times 3.5 = 14$}
\begin{parts}
  \item Airy's and Pratt's hypothesis of Isostasy
  \item Erosional geomorphic features of river in young stage (upper course)
  \item Landscape Plain
  \item Flood and their mitigation
\end{parts}

\medskip\noindent\textbf{Question 5.} Discuss with suitable examples: \pts{$2 \times 7 = 14$}
\begin{parts}
  \item Depositional geomorphic features formed by glacier
  \item Briefly describe the different types of earth movements
\end{parts}

\medskip\noindent\textbf{Question 6.} Discuss the following: \pts{$2 \times 7 = 14$}
\begin{parts}
  \item Types of unconformities and their geological significance
  \item Criteria of recognition of faults in the field and on the maps
\end{parts}

\medskip\noindent\textbf{Question 7.} Differentiate between the following: \pts{$4 \times 3.5 = 14$}
\begin{parts}
  \item Slaty cleavage and strain slip cleavage
  \item Dip fault and dip slip fault
  \item Klippe and Window
  \item Isoclinal and Monoclinal fold
\end{parts}

\bigskip\noindent{\large\bfseries SECTION --- C (Compulsory)}
\rule{\linewidth}{0.4pt}\smallskip

\noindent\textbf{Question 8.} Describe each of the following in two to three sentences. Give suitable diagrams where applicable: \pts{$7 \times 2 = 14$}
\begin{romanparts}
  \item Difference between active remote sensing and passive remote sensing
  \item Distinction within Bouguer anomaly and topographic anomaly
  \item Define the Median Mass
  \item Listric fault
  \item Horst and Graben
  \item Autochthonous and allochthonous beds
  \item Positive and negative flower structure
\end{romanparts}

\vfill\rule{\linewidth}{0.4pt}
\begin{center}\small
\href{https://bhu-exam.vercel.app/}{Click here to visit our website}\\[4pt]
\href{https://me-aryan.vercel.app/}{Click here to visit our contributor}\\[4pt]
\href{https://quashan-me.vercel.app/}{Click here to visit our contributor}
\end{center}
\end{document}
